In the present section\footnote{We use \textit{right} modules in this 
section and read from right to left: we think of the paths 
in the quiver as applying morphisms/functors. The reason for this 
will become clear in Subsection~\ref{sub-ainfty}.} we identify 
the quiver for the blocks $\T_{\lambda}$. We show how the 
quiver description can be used to define a grading on $\T$. We 
obtain a grading on $\Endo(\T)$ as well.



Additionally, we show how the translation functors of $\T_{\lambda}$ 
can be used to get a categorification of the action of the 
braid group on the (split) Grothendieck group of $\T_{\lambda}$ which realizes the 
Burau representation on the (split) Grothendieck group. Moreover, we show how 
the category of projective endofunctors on $\T_{\lambda}$ can be used to 
categorify 
(parts of) the Temperley-Lieb $\bQ(v)$-algebras (where the grading decategorifies to the 
indeterminate $v$).


By an algebra we always mean an associative, but not necessarily unital $\bC$-algebra.


\subsection{The quiver algebras \texorpdfstring{$\Am$}{Qm} and \texorpdfstring{$\Ai$}{Ain}}\label{sub-qalg}
Let $m\in\bN$. We consider the following quiver
\begin{equation}\label{eq-quiver}
\xymatrix{
  \raisebox{0.1cm}{\xy(0,0)*{\bullet};(0,-2.5)*{\scriptstyle \bbm};\endxy}  \ar@<-2pt>[r]_/-0.25em/{d_{m}}  &  \raisebox{0.1cm}{\xy(0,0)*{\bullet};(0,-2.5)*{\scriptstyle \bbmd};\endxy}  \ar@<-2pt>[l]_{u_{m-1}}\ar@<-2pt>[r]_{d_{m-1}} & 
\raisebox{0.1cm}{\xy(0,0)*{\bullet};(0,-2.5)*{\scriptstyle \bbmdd};\endxy}  \ar@<-2pt>[l]_{u_{m-2}}\ar@<-2pt>[r]_{d_{m-2}} &  
  \cdots \ar@<-2pt>[l]_{u_{m-3}}\ar@<-2pt>[r]_/+0.2em/{d_{3}} & \raisebox{0.1cm}{\xy(0,0)*{\bullet};(0,-2.5)*{\scriptstyle \bbt};\endxy}  \ar@<-2pt>[l]_/-0.2em/{u_{2}}\ar@<-2pt>[r]_{d_{2}} & \raisebox{0.1cm}{\xy(0,0)*{\bullet};(0,-2.5)*{\scriptstyle \bbo};\endxy}  \ar@<-2pt>[l]_{u_{1}}\ar@<-2pt>[r]_{d_{1}} & \raisebox{0.1cm}{\xy(0,0)*{\bullet};(0,-2.5)*{\scriptstyle \bbz};\endxy}  \ar@<-2pt>[l]_{u_0}\\
}
\end{equation}
having $m+1$ vertices $\bbz,\bbo,\dots,\bbm$ and $2m$ arrows called \textit{up and down}:
\[
u_i\colon \bbi\to \bbiu,\;i=0,\dots,m-1,\quad\quad d_i\colon \bbi\to \bbid,\;i=1,\dots,m.
\]



The \textit{path algebra} $P_m$ of the quiver 
from~\eqref{eq-quiver} is defined to 
be the $\bC$-algebra whose underlying 
$\bC$-module is the $\bC$-vector 
space spanned by all finite paths 
with multiplication given by composition 
of paths if possible and zero otherwise.


Recall that a $\bZ$\textit{-graded} 
$\bC$-algebra $A$ (or simply \textit{graded}) is a 
$\bC$-algebra $A=\oplus_{i\in\bZ} A^i$ such that $A^iA^j\subset A^{i+j}$ 
for all $i,j\in\bZ$.
The path algebra $P_m$ is a graded 
$\bC$-algebra\footnote{Note that Khovanov and Seidel 
use free $\bZ$-modules instead of $\bC$-vector spaces. 
In order to avoid too many different overlapping notations, 
we only use the $\bC$-vector space version, but it is not a big problem to work over $\bZ$.}, 
that is
\[
P_m=\bigoplus_{k\in\bZ_{\geq 0}}P^k_m,\quad\text{with}\; P^k_m=\{\text{All paths in }P_m\text{ of length }k\},
\]
since we clearly have 
$P^k_m\circ P^{k^{\prime}}_m\subset P^{k+k^{\prime}}_m$. We 
write $l(\cdot)$ for the grading on $P_m$ (the ``length'').





\begin{defn}\label{def-ksquiver1}(\textbf{Khovanov-Seidel's $m$-quiver algebra})
Let $\Am$ denote the quotient algebra obtained from the path algebra 
$P_m$ for the quiver from~\eqref{eq-quiver} by the defining relations
\begin{gather*}
u_{i}\circ u_{i-1}=0=d_i\circ d_{i+1},\;i=1,\dots,m-1,\quad\quad d_{i+1}\circ u_i=u_{i-1}\circ d_{i},\;i=1,\dots,m-1,\\
d_1\circ u_0=0.
\end{gather*}
The latter is called the \textit{dead-end relation}.
Given two paths $p,p^{\prime}\in \Am$ we write $p^{\prime}p$ 
instead of $p^{\prime}\circ p$. The algebra $\Am$ 
inherits the grading $l(\cdot)$ from $P_m$ 
since all the relations are homogeneous. 
We call $\Am$ for short \textit{KS} $m$\textit{-quiver algebra}.
\end{defn}




We denote, by abuse of notation, the path of length $0$ that 
starts and ends at $\bbi$ also by $\bbi$. Note that 
the $\bbi$ are projectors or idempotents, because 
$\bbi\bbi=\bbi$, $p\bbi=p$ if $p$ starts in $\bbi$ 
and $0$ else and $\bbi p=p$ if $p$ ends in $\bbi$ 
and $0$ else. Moreover, they form a complete set of 
pairwise orthogonal idempotents, that 
is $1=\bbz+\bbo+\dots+\bbm$ and $\bbi\bbj=\delta(\bbi,\bbj)$, 
where $1\in \Am$ is the unit. Note that $0\neq \bbz$, 
$1\neq \bbo$ and the $\bbi$'s are not central 
for $m>0$, since e.g. $0=\bbi u_i\neq u_i\bbi=u_i$ for $i=0,\dots,m-1$.



Moreover, we denote for $i=1,\dots,m$ by $\varepsilon_i=u_{i-1}d_i$ 
the loop that starts at $\bbi$ and goes via $d_i$ to $\bbid$ 
and then back via $u_{i-1}$. Note that the relations imply 
that $\varepsilon_i=d_{i+1}u_i$ (if possible, i.e. if $i+1\leq m$). 
Thus, the $\bC$-algebra $\Am$ has a basis given by $\bbi$ 
(for $i=0,\dots,m$) and $u_i$ (for $i=0,\dots,m-1$) and 
$d_i, \varepsilon_i$ (for $i=1,\dots,m$) with $l(\bbi)=0,l(u_i)=l(d_i)=1$ and $l(\varepsilon_i)=2$.


We call an algebra homomorphism $f\colon A\to B$ between 
graded algebras $A$ and $B$ a 
\textit{homomorphism of graded algebras}, if $f(A^k)\subset B^k$ for all $k\in\bZ$.
We have the following.

\begin{lem}\label{lem-inclusion}
There is a sequence of (non-unital!) inclusions of graded algebras
\[
\xymatrix{
  Q_0 \ar@^{(->}[r]^{\iota_0} &  Q_1  \ar@^{(->}[r]^{\iota_1}  & Q_2  \ar@^{(->}[r]^{\iota_2} & \dots,
}
\]
where $\iota_m\colon Q_m\to Q_{m+1}$ is defined by $\bbi,u_i,d_i,\varepsilon_i\mapsto \bbi,u_i,d_i,\varepsilon_i$ for all suitable indices $i$.\makeqed
\end{lem}



\begin{proof}
This follows because $P_m$ includes 
into $P_{m+1}$ and the set of relations for 
$\Am$ is included in the ones for $Q_{m+1}$. 
That these morphisms respect the grading and are injective is immediate.
\end{proof}




This motivates the following definition.
\begin{defn}\label{def-ksquiver2}(\textbf{Khovanov-Seidel's $\infty$-quiver algebra})
Define $\Ai$ to be the inductive limit 
of the sequence of inclusions of graded 
algebras from Lemma~\ref{lem-inclusion}, that is $\Ai=\varinjlim \Am$. 
We call it for short 
\textit{KS} $\infty$\textit{-quiver algebra}.
\end{defn}



Note that $\Ai$ is a graded $\bC$-vector 
space of countable dimension. Moreover, 
$\Ai$ is an algebra with a 
complete set of pairwise orthogonal 
idempotents $\{\bbi\mid i\in\bN\}$, 
but $\Ai$ is a \textit{non-unital} algebra, 
since the unit would have to be an infinite sum of the $\bbi$'s. Such 
an algebra is sometimes called \textit{idempotented}.


\begin{ex}\label{ex-quiver}
$Q_0$ consists just of $\bC$-multiples of 
$\bbz$. The algebra $Q_{m>0}$ can be visualized as
\[
\xymatrix{
  \raisebox{0.1cm}{\xy(0,0)*{\bullet};(0,-2.5)*{\scriptstyle \bbm};\endxy} \ar@(ul,ur)|{\varepsilon_m} \ar@<-2pt>[r]_{d_{m}} & \cdots \ar@<-2pt>[l]_{u_{m-1}}\ar@<-2pt>[r]_/+0.2em/{d_{i+1}} &  \raisebox{0.1cm}{\xy(0,0)*{\bullet};(0,-2.5)*{\scriptstyle \bbi};\endxy} \ar@(ul,ur)|{\varepsilon_i} \ar@<-2pt>[l]_/-0.2em/{u_{i}}\ar@<-2pt>[r]_{d_{i}} & \cdots \ar@<-2pt>[l]_{u_{i-1}}\ar@<-2pt>[r]_/+0.2em/{d_{3}} & \raisebox{0.1cm}{\xy(0,0)*{\bullet};(0,-2.5)*{\scriptstyle \bbt};\endxy} \ar@(ul,ur)|{\varepsilon_2} \ar@<-2pt>[l]_/-0.2em/{u_{2}}\ar@<-2pt>[r]_{d_{2}} & \raisebox{0.1cm}{\xy(0,0)*{\bullet};(0,-2.5)*{\scriptstyle \bbo};\endxy} \ar@(ul,ur)|{\varepsilon_1} \ar@<-2pt>[l]_{u_{1}}\ar@<-2pt>[r]_{d_{1}} & \raisebox{0.1cm}{\xy(0,0)*{\bullet};(0,-2.5)*{\scriptstyle \bbz};\endxy}  \ar@<-2pt>[l]_{u_0}\\
}
\]
The KS $\infty$-quiver algebra $\Ai$ can be visualized as
\[
\xymatrix{
   \cdots\ar@<-2pt>[r]_/+0.2em/{d_{i+1}} &  \raisebox{0.1cm}{\xy(0,0)*{\bullet};(0,-2.5)*{\scriptstyle \bbi};\endxy} \ar@(ul,ur)|{\varepsilon_i} \ar@<-2pt>[l]_/-0.2em/{u_{i}}\ar@<-2pt>[r]_{d_{i}} & \cdots \ar@<-2pt>[l]_{u_{i-1}}\ar@<-2pt>[r]_/+0.2em/{d_{3}} & \raisebox{0.1cm}{\xy(0,0)*{\bullet};(0,-2.5)*{\scriptstyle \bbt};\endxy} \ar@(ul,ur)|{\varepsilon_2} \ar@<-2pt>[l]_/-0.2em/{u_{2}}\ar@<-2pt>[r]_{d_{2}} & \raisebox{0.1cm}{\xy(0,0)*{\bullet};(0,-2.5)*{\scriptstyle \bbo};\endxy} \ar@(ul,ur)|{\varepsilon_1} \ar@<-2pt>[l]_{u_{1}}\ar@<-2pt>[r]_{d_{1}} & \raisebox{0.1cm}{\xy(0,0)*{\bullet};(0,-2.5)*{\scriptstyle \bbz};\endxy}  \ar@<-2pt>[l]_{u_0}\\
}
\]
We point out that $\Ai$ only has one 
asymmetry coming from the dead-end 
relation. Moreover, $\Ai$ is graded by the 
path length $l(\cdot)$ due to the fact 
that $\Ai$ does not contain infinite paths.
\end{ex}



In order to be able to also consider the 
semisimple blocks $\T_{-1}$ and $\T_{l-1}$ of 
our category $\T$, we also introduce another 
quotient of the path algebra $P_m$, denoted 
by $\At$, (and take a limit as above). Since 
the corresponding module categories should be 
semisimple, the quotient $\At$ is rather trivial 
and we call it \textit{KS trivial $m$-quiver algebra}.




\begin{defn}\label{defn-trivial}(\textbf{Khovanov-Seidel's trivial $m$-quiver algebra})
Let $\At$ denote the quotient algebra 
obtained from the path algebra $P_m$ for 
the quiver from~\eqref{eq-quiver} by the defining relations
\[
u_{i}=0=d_{i+1},\;i=0,\dots,m-1.
\]
The algebra $\At$ inherits the grading 
$l(\cdot)$ from $P_m$, but this grading is trivial. 
$\At$ consists only of orthogonal idempotents $\bbi$ for $i=0,\dots, m$. These idempotents 
form a set of pairwise non-isomorphic, central, orthogonal idempotents which shows 
that $\At$ is a semisimple algebra. Moreover, the algebra $\At$ is clearly isomorphic to $\bC\times\dots\times\bC$ (with $m+1$ factors).



As before we define \textit{Khovanov-Seidel's trivial} $\infty$\textit{-quiver algebra} via $\Att=\varinjlim\At$. We note that this can be visualized as
\[
\xymatrix{
   \cdots &  \raisebox{0.1cm}{\xy(0,0)*{\bullet};(0,-2.5)*{\scriptstyle \bbi};\endxy}  & \cdots  & \raisebox{0.1cm}{\xy(0,0)*{\bullet};(0,-2.5)*{\scriptstyle \bbt};\endxy} & \raisebox{0.1cm}{\xy(0,0)*{\bullet};(0,-2.5)*{\scriptstyle \bbo};\endxy} & \raisebox{0.1cm}{\xy(0,0)*{\bullet};(0,-2.5)*{\scriptstyle \bbz};\endxy}\\
}
\]
where we take the inductive limit as above.
\end{defn}



For convenience we set $u_m=d_0=\varepsilon_0=0$ for all $\Am$ and $d_0=\varepsilon_0=0$ for $\Ai$ in the following.



\subsection{Combinatorics of the graded, right \texorpdfstring{$\Am$}{Qm}- and \texorpdfstring{$\Ai$}{Ain}-modules}\label{sub-ammod}
Recall that, if $A$ denotes some 
graded $\bC$-algebra, then a (right) 
$A$-module $M$ is called \textit{$\bZ$-graded} (or simply \textit{graded}), if
\[
M=\bigoplus_{k\in\bZ}M^k\quad\text{and}\quad M^k\cdot A^{k^{\prime}}\subset M^{k+k^{\prime}}.
\]
An $A$-module homomorphism between 
graded modules $f\colon M\to N$ is called 
\textit{degree preserving}, if $f(M^k)\subset N^k$ 
for all $k\in\bZ$ and \textit{homogeneous 
(of degree $d\in\bZ$)} if $f(M^k)\subset N^{k+d}$ for all $k\in\bZ$.
We denote by $\ModgrA$ the category of 
\textit{graded}, finitely generated $A$-modules 
whose morphisms from $M$ to $M^{\prime}$ are given by
\begin{equation}\label{eq-gradhom}
\Hom_{\ModgrA}(M,M^{\prime})=\bigoplus_{s\in\bZ}\Hom_{A}(M,M^{\prime}\langle s\rangle)_{0},
\end{equation}
where the zero should mean \textit{degree preserving} 
morphisms and $\langle s\rangle$ 
denotes a shift as below. Thus, all morphisms in $\ModgrA$ are finite 
direct sums of homogeneous morphisms.




The endofunctor $\cdot\langle s\rangle\colon\ModgrA\to\ModgrA$, 
called \textit{shift\footnote{In our convention a positive number shifts the degree \textit{up}.} 
by $s\in\bZ$} sends the $k-s$-th degree part of a module to the 
$k$-th of the shift $M\langle s\rangle$, that is, $M\langle s\rangle^k=M^{k-s}$.



We use similar notions for left modules 
(we denote such categories by e.g. $\AModgr$) 
or projective modules (we denote such categories 
by e.g. $\pModgrA$). We \textit{only} 
work in categories of (graded) finitely generated (projective) modules.



$\Am$ acts on itself from the right 
by \textit{pre}-composition and from the 
left by \textit{post}-composition of paths.
This makes a difference: 
let us denote by $P_i$ the 
left ideal of $\Am$ generated 
by $\bbi$. Similar, ${}_iP$ denotes 
the right ideal generated by $\bbi$. We have as $\bC$-vector spaces
\[
P_i=\bC\bbi\oplus\bC u_i\oplus\bC d_i\oplus\bC\varepsilon_i,\quad\quad{}_iP=\bC\bbi\oplus\bC u_{i-1}\oplus\bC d_{i+1}\oplus\bC\varepsilon_i,
\]
with homogeneous components of degree $0,1,1,2$ (from left to right).




The $P_i$ and the ${}_iP$ are graded 
$\Am$-modules (left and right respectively) that can be visualized as
\[
P_i=\xymatrix{
  \cdots  & \raisebox{0.1cm}{\xy(0,0)*{\bullet};(0,-2.5)*{\scriptstyle \bbi};\endxy}\ar@(ul,ur)|{\varepsilon_i}  \ar@<-2pt>[l]_/-0.2em/{u_i}\ar@<-2pt>[r]_/-0.2em/{d_i} &\cdots \\
},\quad\quad{}_iP=\xymatrix{
   \cdots \ar@<-2pt>[r]_{d_{i+1}} & \raisebox{0.1cm}{\xy(0,0)*{\bullet};(0,-2.5)*{\scriptstyle \bbi};\endxy}\ar@(ul,ur)|{\varepsilon_i}   &\cdots \ar@<-2pt>[l]_{u_{i-1}}.\\
}
\]
Thus, the $P_i$'s and the ${}_iP$'s 
are projective, since we have (as left, respectively right, $\Am$-modules)
\[
\Am=\bigoplus_{i=0}^m P_i,\quad\quad \Am=\bigoplus_{i=0}^m {}_iP.
\]
Moreover, they are all indecomposable because 
they are the projective covers of the simple left, respectively right, $\Am$-modules $L_i$ 
and ${}_iL$ 
obtained from $P_i$ 
and ${}_iP$ by killing $u_i,d_i$ and $\varepsilon_i$ 
respectively $u_{i-1}$, $d_{i+1}$ and $\varepsilon_i$.
Furthermore, it is easy to see that all 
indecomposable left, respectively right, 
$\Am$-modules are of the form $P_i\langle s\rangle$, 
respectively ${}_iP\langle s\rangle$, for some $s\in\bZ$: 
this follows directly from the fact that $\Am$ is 
finite-dimensional. Thus, all indecomposable projective $\Am$-modules are 
direct summands of $\Am$ considered as a (left or right) $\Am$-module.




\begin{ex}\label{ex-decomp}
One easily checks that the same is true for 
the corresponding notions for $\Ai$ and 
that one gets a decomposition into ${}_iP$'s as 
follows (we have indicated the neighbouring ${}_iP$'s 
using alternating colors and we for simplicity assume that the $i$ below is odd)
\[
\xymatrix{
   \cdots\ar@[blue]@<-2pt>[r]_/+0.2em/{{\color{blue}d_{i+1}}} &  {\color{blue}\raisebox{0.1cm}{\xy(0,0)*{\bullet};(0,-2.5)*{\scriptstyle \bbi};\endxy}} \ar@[blue]@(ul,ur)|{{\color{blue}\varepsilon_i}} \ar@[green]@<-2pt>[l]_/-0.2em/{{\color{green}u_{i}}}\ar@[green]@<-2pt>[r]_{{\color{green}d_{i}}} & \cdots \ar@[blue]@<-2pt>[l]_{{\color{blue}u_{i-1}}}\ar@[green]@<-2pt>[r]_/+0.2em/{{\color{green}d_{3}}} & {\color{green}\raisebox{0.1cm}{\xy(0,0)*{\bullet};(0,-2.5)*{\scriptstyle \bbt};\endxy}} \ar@[green]@(ul,ur)|{{\color{green}\varepsilon_2}} \ar@[blue]@<-2pt>[l]_/-0.2em/{{\color{blue}u_{2}}}\ar@[blue]@<-2pt>[r]_{{\color{blue}d_{2}}} & {\color{blue}\raisebox{0.1cm}{\xy(0,0)*{\bullet};(0,-2.5)*{\scriptstyle \bbo};\endxy}} \ar@[blue]@(ul,ur)|{{\color{blue}\varepsilon_1}} \ar@[green]@<-2pt>[l]_{{\color{green}u_{1}}}\ar@[green]@<-2pt>[r]_{{\color{green}d_{1}}} & {\color{green}\raisebox{0.1cm}{\xy(0,0)*{\bullet};(0,-2.5)*{\scriptstyle \bbz};\endxy}} \ar@[blue]@<-2pt>[l]_{{\color{blue}u_0}}\\
}
\]
We encourage the reader to draw the 
decomposition into $P_i$'s. Again, all 
finitely generated, graded projective, 
indecomposable right $\Ai$-modules are, up 
to a shift, of the form ${}_iP$. Same for 
the left modules and the $P_i$'s. This can 
be deduced directly, but is 
true in more generality for idempotented 
algebras, see for example~\cite[Proposition~5.3.1]{klms}.
\end{ex}



Similarly for $\At$ and $\Att$ which decompose into the 
corresponding (left and right) modules denoted 
by $P_i^{\mathrm{triv}}$ and ${}_iP^{\mathrm{triv}}$.




\subsection{Endofunctors on \texorpdfstring{$\pModgr$-$\Am$}{pMod-Qm} and \texorpdfstring{$\pModgr$-$\Ai$}{pMod-Qin}}\label{sub-endo}
Set\footnote{Note that we use a different convention than Khovanov and Seidel for the gradings by shifting ${}_iP$ down by one.} 
$B_i=P_i\otimes_{\bC}{}_iP\langle -1\rangle$ for all 
indices $i=1,\dots,m$. Note that the $B_i$'s are 
graded $\Am$-bimodules with the tensor product 
$p\otimes p^{\prime}$ of degree $l(p\otimes p^{\prime})=l(p)+l(p^{\prime})$ 
for $p$ of degree $l(p)$ and $p^{\prime}$ of degree $l(p^{\prime})$. 
We define functors
\[
\mathcal U_i\colon \pModgrAm\to \pModgrAm,\;\mathcal U_i= \cdot \otimes_{\Am}B_i\quad\text{for}\quad i=1,\dots,m.
\]
Note that we have
\begin{equation}\label{eq-ufunctors}
{}_{i^{\prime}}P\otimes_{\Am}P_i\cong\biggl\{
\begin{matrix}
 \bC\bbi\oplus\bC\varepsilon_i,\; \text{if }|i-i^{\prime}|=0, & & \bC u_{i},\;\text{if }i^{\prime}-i=1,\\
 \bC d_{i},\;\text{if }i-i^{\prime}=1, & & 0,\;\text{if }|i-i^{\prime}|>1.
\end{matrix}
\end{equation}
This can be easily seen by considering (here $i=i^{\prime}$ as an example)
\[
\xymatrix{
   \cdots \ar@<-2pt>[r]_{d_{i+1}} & \raisebox{0.1cm}{\xy(0,0)*{{\color{green}\bullet}};(0,-2.5)*{{\color{green}\scriptstyle \bbi}};\endxy}\ar@[green]@(ul,ur)|{{\color{green}\varepsilon_i}}   &\cdots \ar@<-2pt>[l]_{u_{i-1}}\\
}\otimes_{\Am}\xymatrix{
  \cdots  & \raisebox{0.1cm}{\xy(0,0)*{{\color{green}\bullet}};(0,-2.5)*{{\color{green}\scriptstyle \bbi}};\endxy}\ar@[green]@(ul,ur)|{{\color{green}\varepsilon_i}}  \ar@<-2pt>[l]_/-0.3em/{u_i}\ar@<-2pt>[r]_/-0.3em/{d_i} &\cdots, \\
}
\]
where we have illustrated the 
overlapping pieces of ${}_iP$ and $P_i$. This clearly implies
\begin{equation}\label{eq-tlcat}
\mathcal U_i ({}_{i^{\prime}}P)\cong\begin{cases}{}_iP\langle -1\rangle\oplus {}_iP\langle +1\rangle,&\text{if }|i-i^{\prime}|=0,\\
{}_iP,&\text{if }|i-i^{\prime}|=1,\\
0,&\text{if }|i-i^{\prime}|>1.\end{cases}
\end{equation}
We note that we can see ${}_{i^{\prime}}P\otimes_{\Am}P_i$ as a 
graded, right $\At$-module where the action of $\bbi$ kills 
everything that does not start in $\bbi$. Note that, in this 
notation, $\bC\bbi,\bQ u_{i},\bC d_{i}$ and $\bC\varepsilon_i$ 
are all one dimensional $\At$-modules concentrated 
in degrees $0,1,1$ and $2$ respectively.
Hence, as graded, 
right $\At$-modules, 
$\bC u_{i}\cong\bC \bbi\langle+1\rangle$, $\bC d_{i}\cong\bC \bbi\langle+1\rangle$ and $\bC \varepsilon_{i}\cong\bC \bbi\langle+2\rangle$.



\begin{cor}\label{cor-tlfunctor}
The functors $\mathcal U_i\colon\pModgrAm\to\pModgrAm$ satisfy the following.
\begin{itemize}
\item[(a)] $\mathcal U_i\circ\mathcal U_i\cong \mathcal U_i\langle -1\rangle\oplus \mathcal U_i\langle +1\rangle$ for $i=1,\dots,m$.
\item[(b)] Let $i^{\prime}=i\pm 1$ such that $i,i^{\prime}\in\{1,\dots,m\}$. Then $\mathcal U_i\circ\mathcal U_{i^{\prime}}\circ\mathcal U_i\cong\mathcal U_i$.
\item[(c)] $\mathcal U_i\circ\mathcal U_{i^{\prime}}\cong 0$, if $|i-i^{\prime}|>1$.\makeqed
\end{itemize}
\end{cor}



\begin{proof}
This follows directly from~\eqref{eq-tlcat}.
\end{proof}



Finally we set
\[
\mathcal U_{\mathrm{even}}=\!\!\bigoplus_{1\leq i, 2i\leq m}\!\!\mathcal U_{2i},\quad\quad\mathcal U_{\mathrm{odd}}=\!\!\bigoplus_{0\leq i, 2i+1\leq m}\!\!\mathcal U_{2i+1}.
\]
By Corollary~\ref{cor-tlfunctor} we see that
\begin{equation}\label{eq-evenodd}
\mathcal U_{\mathrm{even}}\circ \mathcal U_{\mathrm{even}}\cong \mathcal U_{\mathrm{even}}\langle -1\rangle\oplus \mathcal U_{\mathrm{even}}\langle +1\rangle,\quad\quad\mathcal U_{\mathrm{odd}}\circ \mathcal U_{\mathrm{odd}}\cong \mathcal U_{\mathrm{odd}}\langle -1\rangle\oplus \mathcal U_{\mathrm{odd}}\langle +1\rangle.
\end{equation}
It is clear that we can easily adapt 
the definition of the $\mathcal U_i$'s 
to the $\Ai$ case. We denote these 
functors by $\mathcal U^{\infty}_i$ 
for all $i\in\bZ_{>0}$. These functors satisfy completely similar relations as before.



We also define $\mathcal U^{\infty}_{\mathrm{even}}$ 
and $\mathcal U^{\infty}_{\mathrm{odd}}$ completely 
analogously by taking the direct sum over all 
$i\in\bZ_{>0}$. The reader is still invited to 
check that these two satisfy similar relations as in~\eqref{eq-evenodd}.



We also refine 
the definitions of the functors $\mathcal U_i$ 
by \textit{factoring through} the category of 
graded (right) $\At$-modules which we denote by 
$\ModgrAt$. To be precise, we define $\mathcal U_i^t$ and $\mathcal U^i_t$ via
\begin{gather*}
\mathcal U_i^t\colon\ModgrAm\to\ModgrAt,\; \mathcal U_i^t=\cdot\otimes_{\Am}P_i
\\
\mathcal U^i_t\colon\ModgrAt\to\ModgrAm,\; \mathcal U^i_t=\cdot\otimes_{\bC}{}_iP\langle -1\rangle.
\end{gather*}
It is immediate that $\mathcal U_i=\mathcal U^i_t\circ\mathcal U^t_i$. 
We define even and odd versions of these refinements 
as before where we use $\mathcal U^t_{\mathrm{even}}$ 
and $\mathcal U_t^{\mathrm{even}}$ as notations (similar for the odd versions).



The following lemma can be seen, by Remark~\ref{rem-cato} below, 
as an analog of~\cite[Theorem~8.4]{st2}.



\begin{lem}\label{lem-adjoint}
As graded functors: $\mathcal U^t_i\langle+1\rangle$ is 
the right adjoint of $\mathcal U_t^i$, $\mathcal U^t_i$ 
is the left adjoint of $\mathcal U_t^i\langle-1\rangle$. 
Similar for the even/odd versions. The functors $\mathcal U_i$, $\mathcal U_{\mathrm{even}}$ 
and $\mathcal U_{\mathrm{odd}}$ are all graded self adjoint.\makeqed
\end{lem}



\begin{proof}
This is a case-by-case check. 
We use that, as graded $\bC$-vector spaces, we have
\begin{equation}\label{eq-homps}
\Hom_{\pModgrAm}({}_iP,{}_{i^{\prime}}P)\cong\biggl\{
\begin{matrix}
 \bC\bbi\oplus\bC\varepsilon_i,\; \text{if }|i-i^{\prime}|=0, & & \bC u_{i},\;\text{if }i^{\prime}-i=1,\\
 \bC d_{i},\;\text{if }i-i^{\prime}=1, & & 0,\;\text{if }|i-i^{\prime}|>1.
\end{matrix}
\end{equation}
The isomorphism is given by mapping a 
path $p$ on the right-hand side to the 
homomorphism ${}_iP\to {}_{i^{\prime}}P$ 
given by left multiplication 
(\textit{post}-composition) with $p$. 
Note that this is a homogeneous morphism of degree $l(p)$.


We only 
prove that $\mathcal U^t_i\langle+1\rangle$ is 
the right adjoint of $\mathcal U_t^i$ and leave the rest to the reader.
To this end, we have to show that there are isomorphisms of graded vector spaces
\[
\Hom_{\pModgrAm}(\mathcal U_t^i(\bC\bbi),{}_{i^{\prime}}P)\cong\Hom_{\pModgrAt}(\bC\bbi,\mathcal U^t_i\langle+1\rangle({}_{i^{\prime}}P)),
\]
where we can restrict to check only for objects 
as above by additivity. Note that we have three 
cases depending on $i-i^{\prime}$, since the case 
$|i-i^{\prime}|>1$ is clear.


$\boldsymbol{i=i^{\prime}}$. The left-hand side is 
$\Hom_{\pModgrAm}({}_iP\langle-1\rangle,{}_iP)\cong \bC\bbi\langle+1\rangle\oplus \bC\varepsilon_i\langle+1\rangle$, which follows from~\eqref{eq-homps}. 
Moreover, by~\eqref{eq-ufunctors} and the shift by $+1$, we get the same for the right side.



$\boldsymbol{i^{\prime}=i+1}$. The left-hand 
side gives $\Hom_{\pModgrAm}({}_iP\langle-1\rangle,{}_{i+1}P)\cong \bC u_i\langle+1\rangle$. This follows again by~\eqref{eq-homps}. The right-hand side gives, 
by~\eqref{eq-ufunctors} and the shift by 
$+1$, the same result.

$\boldsymbol{i^{\prime}=i-1}$. Similarly to the case $\boldsymbol{i^{\prime}=i+1}$.


Thus, we have an adjoint pair 
$(\mathcal U_t^i,\mathcal U^t_i\langle+1\rangle)$. 
The other statements follow analogously.
\end{proof}


The same holds, mutatis mutandis, for the infinity 
versions as well.


%We again have no problems to define the $\infty$-versions. 
%We denote them by $\mathcal U_i^{\infty}$ and 
%$\mathcal U^i_{\infty}$. The even versions of 
%these are denoted by $\mathcal U_i^{\mathrm{even}\cdot\infty}$ 
%and $\mathcal U^i_{\mathrm{even}\cdot\infty}$ and similar for the odd versions.



\subsection{\texorpdfstring{$\Ai$}{Ain} and the tilting category \texorpdfstring{$\T$}{T}}\label{sub-ainfty}
Recall that the category 
$\T_{\lambda}$ denotes the 
$\lambda$-block of $\T$ for $\lambda$ 
in the fundamental alcove $\mathcal A_0$, see Subsection~\ref{sub-blocks}. 
We fix any such $\lambda$ and denote everything 
using the simplified notation without the $\lambda$'s. In particular, recall our notation from Definition~\ref{defn-tiltgen} for the tilting generators $\Tn{\bbin}$ and $\Tn{\leq\bbm}$.
We denote by $\T_{\lambda}(\leq\bbm)$ the 
full subcategory of $\T_{\lambda}$ consisting of 
objects whose indecomposable summands are all from 
the set $\{\Tn{\lambda_0},\dots,\Tn{\lambda_m}\}$ for some fixed $m\in\bN$.
We also use completely similar notations for 
the walls $\mu$ by indicating this case with a superscript $t$. 
%Note that (similarly for $\Tn{\bbin}^t$)
%\[
%\Tn{\leq\bbm}^t\cong\begin{cases}\bigoplus_{i=0}^{\lfloor m/2\rfloor} \Tn{\mu_{2i}}, &\text{if }\mu=-1,\\ \bigoplus_{i=0}^{\lfloor m/2\rfloor} \Tn{\mu_{2i+1}}, &\text{if }\mu=l-1.\end{cases}
%\]




\begin{prop}\label{prop-isoofalg1}
We have
\[
\End_{\Uq}(\Tn{\leq\bbm})\cong \Am\quad\text{and}\quad\End_{\Uq}(\Tn{\leq\bbm}^t)\cong \At,
\]
which are isomorphisms of $\bC$-algebras.\makeqed
\end{prop}



\begin{proof}
We are going to construct these 
isomorphisms explicitly. Moreover, 
if we simply write $\bbi$ for $\Tn{\bbi}$, 
then we can visualize $\End_{\Uq}(\Tn{\leq\bbm})$ 
using Proposition~\ref{prop-link} for all $m>0$ as
\[
\xymatrix{
  \raisebox{0.1cm}{\xy(0,0)*{\bullet};(0,-2.5)*{\scriptstyle \bbm};\endxy} \ar@(ul,ur)|{\varepsilon_m} \ar@<-2pt>[r]_{d_{m}} & \cdots \ar@<-2pt>[l]_{u_{m-1}}\ar@<-2pt>[r]_/+0.2em/{d_{i+1}} &  \raisebox{0.1cm}{\xy(0,0)*{\bullet};(0,-2.5)*{\scriptstyle \bbi};\endxy} \ar@(ul,ur)|{\varepsilon_i} \ar@<-2pt>[l]_/-0.2em/{u_{i}}\ar@<-2pt>[r]_{d_{i}} & \cdots \ar@<-2pt>[l]_{u_{i-1}}\ar@<-2pt>[r]_/+0.2em/{d_{3}} & \raisebox{0.1cm}{\xy(0,0)*{\bullet};(0,-2.5)*{\scriptstyle \bbt};\endxy} \ar@(ul,ur)|{\varepsilon_2} \ar@<-2pt>[l]_/-0.2em/{u_{2}}\ar@<-2pt>[r]_{d_{2}} & \raisebox{0.1cm}{\xy(0,0)*{\bullet};(0,-2.5)*{\scriptstyle \bbo};\endxy} \ar@(ul,ur)|{\varepsilon_1} \ar@<-2pt>[l]_{u_{1}}\ar@<-2pt>[r]_{d_{1}} & \raisebox{0.1cm}{\xy(0,0)*{\bullet};(0,-2.5)*{\scriptstyle \bbz};\endxy}  \ar@<-2pt>[l]_{u_0}\\
}
\]
Thus, we consider the 
(by Proposition~\ref{prop-link} well-defined) $\bC$-algebra homomorphism
\[
\phi\colon \Am\to\End_{\Uq}(\Tn{\leq\bbm}), u_i\mapsto u_i, d_i\mapsto d_i,\varepsilon_i\mapsto \varepsilon_i
\]
for all suitable indices $i$. 
Using Corollary~\ref{cor-homspaces} 
we see that this is an isomorphism. 
The case $m=0$ and the 
``trivial case'' cases are immediate which proves the statement.
\end{proof}



The morphisms of $\End_{\Uq}(\Tn{\bbin})$ are 
``matrices'' $f=\prod_{i,i^{\prime}\in\bN} f_{i^{\prime}i}$ 
with $f_{i^{\prime}i}\colon \Tn{\bbi}\to\Tn{\bbi^{\prime}}$ and each 
row consists only of finitely many non-zero entries 
(this follows from Corollary~\ref{cor-homspaces}). 
Now $\End^{\mathrm{fs}}_{\Uq}(\Tn{\bbin})$ should be the 
(non-unital) subalgebra of $\End_{\Uq}(\Tn{\bbin})$ 
of all matrices with only finitely many non-zero rows, i.e.
\[
\End^{\mathrm{fs}}_{\Uq}(\Tn{\bbin})\cong \bigoplus_{i,i^{\prime}=0}^{\infty} \Hom_{\Uq}(\Tn{\bbi},\Tn{\bbi^{\prime}}).
\]
Similarly, of course, in the ``trivial'' case.



\begin{prop}\label{prop-isoofalg2}
We have
\begin{gather*}
\begin{aligned}
\End^{\mathrm{fs}}_{\Uq}(\Tn{\bbin})&\cong\bigoplus_{i,i^{\prime}=0}^{\infty}\Hom_{\Uq}(\Tn{\bbi},\Tn{\bbi^{\prime}})\cong \Ai,
\\
\End^{\mathrm{fs}}_{\Uq}(\Tn{\bbin}^t)&\cong\bigoplus_{i,i^{\prime}=0}^{\infty}\Hom_{\Uq}(\Tn{\bbi}^t,\Tn{\bbi^{\prime}}^t)\cong \Att,
\end{aligned}
\end{gather*}
which are isomorphisms of $\bC$-algebras.\makeqed
\end{prop}



\begin{proof}
We only show the first isomorphism since the other is clear.
By using Corollary~\ref{cor-homspaces} and the definition 
of $\End^{\mathrm{fs}}_{\Uq}(\Tn{\bbin})$ from above we 
see that, for each $f\in\End^{\mathrm{fs}}_{\Uq}(\Tn{\bbin})$, 
there exists some $m\in\bN$ such that $f_{i^{\prime}i}=0$ for $i,i^{\prime}>m$.
Those matrices describe elements in $\End_{\Uq}(\Tn{\leq\bbm})$ by ``forgetting'' the parts for $i,i^{\prime}>m$. Thus, we have
\[
\End^{\mathrm{fs}}_{\Uq}(\Tn{\bbin})\cong \varinjlim \End_{\Uq}(\Tn{\leq\bbm}).
\]
By Proposition~\ref{prop-isoofalg1} we see that this implies
\[
\End^{\mathrm{fs}}_{\Uq}(\Tn{\bbin})\cong \varinjlim \Am\cong\Ai.
\]
This finishes the proof.
\end{proof}



We are now able to relate 
$\T_{\lambda}(\leq \bbm)$ to 
$\pModAm$ and thus, $\T_{\lambda}$ to $\pModAi$. 
We define
\begin{gather*}
\bV\colon\T_{\lambda}(\leq \bbm)\to\ModAm, \bV(M)=\Hom_{\Uq}(\Tn{\leq\bbm},M),
\\
\bVi\colon\T_{\lambda}\to\ModAi, \bVi(M)=\Hom^{\mathrm{fs}}_{\Uq}(\Tn{\bbin},M)\cong \bigoplus_{i=0}^{\infty}\Hom_{\Uq}(\Tn{\bbi},M),
\end{gather*}
for $M$ in either $\Ob(\T(\leq \bbm))$ or $\Ob(\T_{\lambda})$ 
respectively where the isomorphism follows from 
Proposition~\ref{prop-isoofalg2}. Similarly for 
the walls $\mu$ where we indicate the functors 
$\bV^t$ and $\bVi^t$ (and everything else) with 
superscripts $t$. We note that the hom-spaces above 
are right $\Am$-modules (or $\Ai$, $\At$ or $\Ait$-modules) by \textit{pre}-composition.



Note that $\Theta_s,\Theta_t\in\Ob(\End(\T_{\lambda}))$ 
from Definition~\ref{defn-trans} clearly descent 
to $\T_{\lambda}(\leq\bbm)$ for all $m\in\bN$. 
We denote these restrictions, by abuse of notation, 
also by $\Theta_s,\Theta_t$. Furthermore, 
recall the notation $\Theta_{s\text{ or }t}^i$ for $i\in\bZ_{>0}$ 
which are defined via ($p_{\bbi}$ projects to the $\Tn{\bbi}$-part)
\[
\Theta_{s\text{ or }t}^i=\begin{cases}p_{\bbi}\circ\Theta_{s}, &\text{if }i\text{ is even },\\
p_{\bbi}\circ\Theta_{t}, &\text{if }i\text{ is odd }.
\end{cases}
\]
We use the 
same conventions and notations in the finite case by killing everything above $m$.



Similarly: the functors $\mathcal T_{\lambda}^{\mu}$ 
and $\mathcal T^{\lambda}_{\mu}$ from Definition~\ref{defn-onout} also 
descent to the finite case where we again denote them by the same symbols.



We consider the right side for the tilting embedding theorem 
(finite and infinite) as ungraded modules/functors. 
Moreover $\pMod$ always denotes categories of projective modules.



\begin{thm}(\textbf{Tilting embedding theorem - finite version})\label{thm-struktursatz1}
We have the following.
\begin{itemize}
\item[(a)] The functor $\bV$ is fully faithful.
\item[(b)] The functor $\bV$ sends objects 
$M\in\Ob(\T_{\lambda}(\leq\bbm))$ to projective 
$\Am$-modules and is an equivalence  between $\T_{\lambda}(\leq\bbm)$ and $\pModAm$.
\item[(c)] We have the following commuting diagram.
\[
\begin{xy}
  \xymatrix{
      \T_{\lambda}(\leq\bbm) \ar[r]^/-0.1em/{\bV} \ar[d]_{\Theta^i_{s\text{ or }t}}    &   \pModAm \ar[d]^{\mathcal U_{i}} \\
      \T_{\lambda}(\leq\bbm) \ar[r]_/-0.1em/{\bV}             &   \pModAm.  
  }
\end{xy}
\]
\item[(d)] We have the following commuting diagrams.
\[
\begin{xy}
  \xymatrix{
      \T_{\lambda}(\leq\bbm) \ar[r]^/-0.1em/{\bV} \ar[d]_{\Theta_s}    &   \pModAm \ar[d]^{\mathcal U_{\mathrm{even}}}  \\
      \T_{\lambda}(\leq\bbm) \ar[r]_/-0.1em/{\bV}             &   \pModAm   
  }
\end{xy}
\quad\raisebox{-0.75cm}{\text{and}}\quad
\begin{xy}
  \xymatrix{
      \T_{\lambda}(\leq\bbm) \ar[r]^/-0.1em/{\bV} \ar[d]_{\Theta_t}    &   \pModAm \ar[d]^{\mathcal U_{\mathrm{odd}}}  \\
      \T_{\lambda}(\leq\bbm) \ar[r]_/-0.1em/{\bV}             &   \pModAm.  
  }
\end{xy}
\]
\item[(e)] The same as in (a)+(b) works for the walls $\mu$ as well by using $\bV^t$.
\item[(f)] We have the following commuting diagrams.
\[
\begin{xy}
  \xymatrix{
      \T_{\lambda}(\leq\bbm) \ar[r]^/-0.1em/{\bV} \ar[d]_{\mathcal{T}_{\lambda}^{-1}}    &   \pModAm \ar[d]^{\mathcal U_{\mathrm{even}}^t}  \\
      \T_{-1}(\leq\bbm)\ar[r]_/-0.1em/{\bV^t}             &   \pModAt  
  }
\end{xy}
\quad\raisebox{-0.75cm}{\text{and}}\quad
\begin{xy}
  \xymatrix{
      \T_{\lambda}(\leq\bbm) \ar[r]^/-0.1em/{\bV} \ar[d]_{\mathcal{T}_{\lambda}^{l-1}}    &   \pModAm \ar[d]^{\mathcal U_{\mathrm{odd}}^t}  \\
      \T_{l-1}(\leq\bbm)\ar[r]_/-0.1em/{\bV^t}             &   \pModAt.  
  }
\end{xy}
\]
\item[(g)] We have the following commuting diagrams.
\[
\begin{xy}
  \xymatrix{
      \T_{-1}(\leq\bbm) \ar[r]^/-0.1em/{\bV^t} \ar[d]_{\mathcal{T}^{\lambda}_{-1}}    &   \pModAt \ar[d]^{\mathcal U^{\mathrm{even}}_t}  \\
      \T_{\lambda}(\leq\bbm) \ar[r]_/-0.1em/{\bV}             &   \pModAm  
  }
\end{xy}
\quad\raisebox{-0.75cm}{\text{and}}\quad
\begin{xy}
  \xymatrix{
      \T_{l-1}(\leq\bbm) \ar[r]^/-0.1em/{\bV^t} \ar[d]_{\mathcal{T}^{\lambda}_{l-1}}    &   \pModAt \ar[d]^{\mathcal U^{\mathrm{odd}}_t}  \\
      \T_{\lambda}(\leq\bbm) \ar[r]_/-0.1em/{\bV}             &   \pModAm.  
  }
\end{xy}
\hspace{1.95cm}
\raisebox{-1.55cm}{\makeqed}
\hspace{-1.95cm} 
\]
\end{itemize}
\end{thm}



\begin{proof}
We are careless with the indices in the proof 
and hope that this does not cause confusion. All 
appearing indices should be such that everything 
below $0$ or above $m$ is defined to be zero.



\textbf{(a).} By Proposition~\ref{prop-tilt} part 
(c) and Lemma~\ref{lem-tiltcat} part (a) it is enough 
to verify the statement for $M=\Tn\bbi$ for $i=0,\dots,m$. 
Moreover, by using Corollary~\ref{cor-homspaces} together 
with the discussion in Subsection~\ref{sub-ammod}, we 
get isomorphisms $\bV(\Tn\bbi)=\Hom_{\Uq}(\Tn{\leq \bbm},\Tn\bbi)\cong {}_iP$. 
Thus, $\bV$ is fully faithful.



\textbf{(b).} From (a) by using the discussion in 
Subsection~\ref{sub-ammod}.



\textbf{(c).} This follows by combining 
$\bV(\Tn\bbi)=\Hom_{\Uq}(\Tn{\leq \bbm},\Tn\bbi)\cong {}_iP$ 
with~\eqref{eq-tlcat} and part (c) of Proposition~\ref{prop-functors}.



\textbf{(d).} This follows directly from (c) 
and the definition of $\mathcal U_{\mathrm{even}}$ and $\mathcal U_{\mathrm{odd}}$. 



\textbf{(e).} The ``trivial'' cases are clear.



\textbf{(f).} We do the left case here. Note that $\bV(\Tn{\bbi})\cong {}_iP$. 
Thus, using~\eqref{eq-ufunctors}, we see that 
$\mathcal U^{\mathrm{even}}_t\bV^t$ sends 
$\Tn{\bbi}$ to $\bC\bbi\oplus\bC\varepsilon_i$ 
for even $i$ and $\bC u_{i-1}\oplus \bC d_{i+1}$ for odd $i$.
This, by the discussion in Subsection~\ref{sub-endo}, 
is isomorphic as right $\At$-modules to 
$\bC\bbi\oplus\bC\bbi$ and $\bC\bbid\oplus\bC\bbiu$ 
respectively. In addition, $\mathcal T^{-1}_{\lambda}$ 
sends $\Tn{\bbi}$ to $\Tn{\mu_i}\oplus\Tn{\mu_i}$ 
(for $i$ even) or to $\Tn{\mu_{i-1}}\oplus\Tn{\mu_{i+1}}$ 
(for $i$ odd). These are then sent to $\bC\bbi\oplus\bC\bbi$ (for $i$ even) 
respectively $\bC\bbid\oplus\bC\bbiu$ (for $i$ odd).



\textbf{(g)}: Analogously to (f).
\end{proof}



\begin{thm}(\textbf{Tilting embedding theorem - infinite version})\label{thm-struktursatz2}
We have the following.
\begin{itemize}
\item[(a)] The functor $\bVi$ is fully faithful.
\item[(b)] The functor $\bVi$ sends objects $M\in\Ob(\T)$ to projective $\Ai$-modules and is an equivalence  between $\T_{\lambda}$ and $\pModAi$.
\item[(c)] We have the following commuting diagram.
\[
\begin{xy}
  \xymatrix{
      \T_{\lambda} \ar[r]^/-0.75em/{\bVi} \ar[d]_{\Theta^i_{s\text{ or }t}}    &   \pModAi \ar[d]^{\mathcal U^{\infty}_{i}}  \\
      \T_{\lambda} \ar[r]_/-0.75em/{\bVi}             &   \pModAi.   
  }
\end{xy}
\]
\item[(d)] We have the following commuting diagrams.
\[
\begin{xy}
  \xymatrix{
      \T_{\lambda} \ar[r]^/-0.75em/{\bVi} \ar[d]_{\Theta_s}    &   \pModAi \ar[d]^{\mathcal U^{\infty}_{\mathrm{even}}}  \\
      \T_{\lambda} \ar[r]_/-0.75em/{\bVi}             &   \pModAi   
  }
\end{xy}
\quad\raisebox{-0.75cm}{\text{and}}\quad
\begin{xy}
  \xymatrix{
      \T_{\lambda} \ar[r]^/-0.75em/{\bVi} \ar[d]_{\Theta_t}    &   \pModAi \ar[d]^{\mathcal U^{\infty}_{\mathrm{odd}}}  \\
      \T_{\lambda} \ar[r]_/-0.75em/{\bVi}             &   \pModAi.  
  }
\end{xy}
\]
\item[(e)] The same as in (a)+(b) works for the walls $\mu$ as well by using $\bVi^t$.
\item[(f)] We have the following commuting diagrams.
\[
\begin{xy}
  \xymatrix{
      \T_{\lambda}(\bbin) \ar[r]^/-0.1em/{\bVi} \ar[d]_{\mathcal{T}_{\lambda}^{-1}}    &   \pModAi \ar[d]^{\mathcal U_{\mathrm{even}}^t}  \\
      \T_{-1}(\bbin) \ar[r]_/-0.1em/{\bVi^t}             &   \pModAtt
  }
\end{xy}
\quad\raisebox{-0.75cm}{\text{and}}\quad
\begin{xy}
  \xymatrix{
      \T_{\lambda}(\bbin) \ar[r]^/-0.1em/{\bVi} \ar[d]_{\mathcal{T}_{\lambda}^{l-1}}    &   \pModAi \ar[d]^{\mathcal U_{\mathrm{odd}}^t}  \\
      \T_{l-1}(\bbin) \ar[r]_/-0.1em/{\bVi^t}             &   \pModAtt.  
  }
\end{xy}
\]
\item[(g)] We have the following commuting diagrams.
\[
\begin{xy}
  \xymatrix{
      \T_{-1}(\bbin) \ar[r]^/-0.1em/{\bVi^t} \ar[d]_{\mathcal{T}^{\lambda}_{-1}}    &   \pModAtt \ar[d]^{\mathcal U^{\mathrm{even}}_t}  \\
      \T_{\lambda}(\bbin) \ar[r]_/-0.1em/{\bVi}             &   \pModAi  
  }
\end{xy}
\quad\raisebox{-0.75cm}{\text{and}}\quad
\begin{xy}
  \xymatrix{
      \T_{l-1}(\bbin) \ar[r]^/-0.1em/{\bVi^t} \ar[d]_{\mathcal{T}^{\lambda}_{l-1}}    &   \pModAtt \ar[d]^{\mathcal U^{\mathrm{odd}}_t}  \\
      \T_{\lambda}(\bbin) \ar[r]_/-0.1em/{\bVi}             &   \pModAi.  
  }
\end{xy}
\hspace{2.3cm}
\raisebox{-1.55cm}{\makeqed}
\hspace{-2.3cm}
\]
\end{itemize}
\end{thm}


\begin{proof}
The proofs of (a)-(g) is similar to 
the proofs given in Theorem~\ref{thm-struktursatz1} 
by using arguments as in Proposition~\ref{prop-isoofalg2}. We leave it to the reader.
\end{proof}



\begin{rem}\label{rem-othermaps}
We encourage the reader to work out the 
translations onto and out of the wall 
commuting diagrams similar to part (c) of the theorems above.
\end{rem}




We note the following corollary for completeness.
\begin{cor}\label{cor-struktursatz}
For all $\lambda,\lambda^{\prime}\in\mathcal A_0$ we have $\T_{\lambda}\cong \T_{\lambda^{\prime}}$. Thus,
\[
\T(\leq\bbm)\cong \bigoplus_{i=0}^{l-2}\pModAm\oplus\bigoplus_{i=0}^{1}\pModAt,
\quad\T\cong \bigoplus_{i=0}^{l-2}\pModAi\oplus\bigoplus_{i=0}^{1}\pModAtt,
\]
with semisimple categories $\pModAt$ and $\pModAtt$.\makeqed
\end{cor}



\begin{proof}
This follows from Lemma~\ref{lem-block} 
and Theorems~\ref{thm-struktursatz1} 
and~\ref{thm-struktursatz2} because the 
right-hand sides are always the same independent of $\lambda\in\mathcal A_0$.
\end{proof}



\begin{ex}(\textbf{How to compute the images of the translation functors on $\Uq$-morphisms})\label{ex-howtofunctor}
Combining the commuting diagrams of the 
Theorems~\ref{thm-struktursatz1} and~\ref{thm-struktursatz2} 
we obtain an explicit way to compute $\Theta_s$ and $\Theta_t$ 
not just on the $\Tn{\bbi}$'s (as we already know by part (c) 
of Proposition~\ref{prop-functors}), but also 
on the $\Uq$-intertwiners in $\Hom_{\Uq}(\Tn{\bbi},\Tn{\bbi^{\prime}})$.



This follows by using~\eqref{eq-homps} together 
with $\Theta_s\leftrightsquigarrow \mathcal U_{\mathrm{even}}$ 
and $\Theta_t\leftrightsquigarrow \mathcal U_{\mathrm{odd}}$. 
For example, given $u_0\colon\Tn{\bbz}\to\Tn{\bbo}$, we can 
compute $\Theta_t(u_0)$ as follows. First note that 
$\bVi\Theta_t(\Tn{\bbz})\cong {}_1P$ and $\bVi\Theta_t(\Tn{\bbo})\cong {}_1P\oplus {}_1P$ by~\eqref{eq-tlcat} and Theorem~\ref{thm-struktursatz2}. By using~\eqref{eq-homps} 
we see
\[
\Theta_t(u_0)=\begin{pmatrix}
 0 \\
 \mathrm{id}
\end{pmatrix}\colon \Tn{\bbo}\to\Tn{\bbo}\oplus \Tn{\bbo}.
\]
We see that this is actually
\[
\Theta_t(u_0)=\begin{pmatrix}
 0 \\
 \mathrm{id} 
\end{pmatrix}\colon \Tn{\bbo}\to\Tn{\bbo}\langle-1\rangle\oplus \Tn{\bbo}\langle+1\rangle,
\]
by using the results we are going to explain in Subsection~\ref{sub-grad}. Indeed, 
the grading is very helpful to calculate the images of the $\Uq$-interwiners, e.g. 
$\Theta_t(u_0)$ can only be zero on the first component, since we will show 
that there will not be  suitable $\Uq$-intertwiners lowering the degree.
\end{ex}



\begin{rem}\label{rem-cato}
Let $\mathfrak{g}=\mathfrak{sl}_{m+1}$ and 
$\mathfrak{p}=\mathfrak{p}_1$ be a maximal 
parabolic subalgebra for $S_1\times S_{m-1}$ 
(for a definition see for example~\cite[Chapter~9]{hum}). 
Denote by $\mathcal O^{\mathfrak{p}}$ the corresponding 
parabolic category $\mathcal O$ (for a definition and 
some properties see e.g.~\cite[Section~9.3]{hum}). Moreover, 
denote by $\mathcal O^{\mathfrak{p}}_0$ the principle block. 
As Khovanov and Seidel show in~\cite[Proposition~2.9]{ks1} 
there is an equivalence of categories 
$\mathcal O^{\mathfrak{p}}_0\cong \ModAm$. 
Thus, by Theorems~\ref{thm-struktursatz1} 
and~\ref{thm-struktursatz2}, we see that 
$\T_{\lambda}$ governs $\mathcal O^{\mathfrak{p}}_0$ for all $m$.
Similarly for the projective endofunctors: 
as explained in~\cite[below Proposition~2.9]{ks1}, 
the functors $\mathcal U_i$ correspond to the translation 
through the $i$-th wall (note that, under the equivalence 
above, the $\mathfrak{sl}_{m+1}$ has Weyl group generators 
$s_1,\dots,s_m$ and therefore translation functors indexed 
by $i=1,\dots,m$). Moreover, $\mathcal U_i^t$ corresponds 
to the translation onto and $\mathcal U^i_t$ corresponds to 
the translation out off the wall functors.
Hence, by Theorems~\ref{thm-struktursatz1} 
and~\ref{thm-struktursatz2}, we can say 
that these are governed by the combinatorics 
of $\mathcal T_{\lambda}^{\mu}$, 
$\mathcal T^{\lambda}_{\mu}$ 
and $\Theta^{\lambda}_{s\text{ or }t}$ from our tilting case.



Another point is, as Khovanov and Seidel explain 
after~\cite[Proposition~2.9]{ks1}, that the path length grading gives rise 
to a \textit{Koszul grading} in the sense of~\cite{bgs} on 
the $\Am$'s, see also~\cite[Example~1.1]{st1}. 
The same holds for $\Ai$. This can either be 
seen ``by hand'' or by using a more general 
theory for (quotients of) not necessary 
finite quiver algebras that can be found for example in~\cite{mos}.
\end{rem}




\begin{rem}\label{rem-webalg}
Khovanov and Seidel's quiver is also 
related to Khovanov's arc algebra 
$H^m$ (for $m$ even) that he introduced 
in~\cite{kh4} to give an algebraic interpretation 
of Khovanov homology. The arc algebra categorifies the 
invariant tensors of $V^{\otimes m}$ where $V$ is the 
two dimensional vector representation of $\Uv$.
In addition, together with Chen he extended in~\cite{chkh} 
this categorification to the full tensor 
product $V^{\otimes m}$ by defining a certain 
subquotient $\prod_k A^{k,m-k}$ (this quotient can 
be seen as the quasi-hereditary cover of $H^m$). 
All of these have a topological interpretation 
as algebras consisting of cobordisms, 
see e.g.~\cite[Section~2]{chkh}, and they are 
graded by \textit{the Euler characteristic} of these cobordisms.
As explained in~\cite[Section~3]{chkh}, 
the $A^{1,m-1}$-part of this subquotient is graded 
isomorphic to $\Am$. Thus, we see that $\T$ governs $A^{1,m-1}$ for all $m$.
Moreover, this gives a hint for a  
generalization of our work: we think 
that analogs of the Theorems~\ref{thm-struktursatz1} and~\ref{thm-struktursatz2} 
can be proven for certain subquotients of the $\mathfrak{sl}_n$ 
generalizations of $H^m$ studied in e.g.~\cite{mack1},~\cite{mpt},~\cite{tub3} or~\cite{tub4}.
\end{rem}




\subsection{The tilting category \texorpdfstring{$\T$}{T} and its graded counterpart \texorpdfstring{$\Tgr$}{Tgr}}\label{sub-grad}
In this subsection we discuss the 
consequences of Theorems~\ref{thm-struktursatz1} 
and~\ref{thm-struktursatz2} with respect to the 
question how the path length grading $l(\cdot)$ 
(see Definition~\ref{def-ksquiver1}) of 
Khovanov-Seidel's $\infty$-quiver algebra relates 
to the tilting category $\T$. We point out that this is non-trivial since Theorems~\ref{thm-struktursatz1} and~\ref{thm-struktursatz2} only say that 
$\T_{\lambda}(\leq \bbm)$ and $\T_{\lambda}$ are 
isomorphic to subcategories of $\pModAm$ and $\pModAi$.



Recall that, for every graded algebra $A$, there is a 
\textit{forgetful functor} 
$\mathrm{forget}\colon\ModgrA\to\ModA$ 
that forgets 
the grading. Using this 
functor, we say that an $A$-module 
$M\in\Ob(\ModA)$ is \textit{gradable} if there 
exists an $A$-module $\tilde M\in\Ob(\ModgrA)$ 
such that $\mathrm{forget}(\tilde M)=M$. Note 
that, by abuse of notation, we usually do not distinguish between $\tilde M$ and $M$.



Recall that $\Hom_{\ModA}(M,M^{\prime})$ 
consists of all (right) $A$-module 
homomorphisms $f\colon M\to M^{\prime}$ 
and $\Hom_{\ModgrA}(\tilde M,\tilde M^{\prime})$ 
consists of (right) $A$-module homomorphisms 
$\tilde f\colon \tilde M\to \tilde M^{\prime}$ as 
in~\eqref{eq-gradhom}. Hence, in the same vein as 
above, we can call an $A$-module homomorphism 
$f\colon M\to M^{\prime}$ \textit{gradable}, 
if there exists $\tilde f\colon\tilde M\to \tilde M^{\prime}$ 
such that $\mathrm{forget}(\tilde f)=f$.
Consequently, we define the following.






\begin{defn}\label{defn-gradable}
We call a (left) $\Uq$-tilting module 
$M\in\Ob(\T)$ \textit{gradable}, if $\bVi(M)$ is 
gradable viewed as a (right) $\Ai$-module. We call a 
$\Uq$-intertwiner $f\colon M\to M^{\prime}\in\Hom_{\T}(M,M^{\prime})$ 
\textit{gradable}, if $\bVi(f)\colon \bVi(M)\to \bVi(M^{\prime})$ is 
a gradable $\Ai$-module homomorphism.



We denote by $\Tgr$ the subcategory of $\T$ consisting 
of gradable $\Uq$-tilting modules and gradable 
$\Uq$-intertwiners. We use similar notations and 
conventions for $\leq\bbm$ and for fixed $\lambda$ or $\mu$.
\end{defn}



The question concerning which objects $M\in\Ob(\T)$ 
are gradable turns out to be surprisingly 
simple in our case. We note that, with respect to 
the rather complicated situation for 
graded category $\mathcal O_0^{\mathfrak{p}}$ 
discussed in e.g.~\cite[Theorem~4.1]{st2}, 
the following proposition is quite remarkable.



\begin{prop}\label{prop-grading}
All objects $M\in\Ob(\T_{\lambda}(\leq\bbm))$, 
$M\in\Ob(\T_{\lambda})$ and $M\in\Ob(\T)$ are gradable. The grading of 
each indecomposable module $\Tn{\bbi}$ is unique up to shifts.\makeqed
\end{prop}



\begin{proof}
We only prove the $\infty$ case and leave 
the other case to the reader. We start by 
considering $\lambda\in\mathcal A_0$ and 
discuss the semisimple case on the walls afterwards.




It is easy to check that, for any graded algebra 
$A$ and any two gradable $A$-modules $M$ and 
$M^{\prime}$, the direct sum $M\oplus M^{\prime}$ 
is also gradable. Thus, by Lemma~\ref{lem-tiltcat} 
part (a), it is enough to consider only the 
$\Tn{\bbi}$'s. As in the proof of the part (b) 
of Theorem~\ref{thm-struktursatz2} 
(or Theorem~\ref{thm-struktursatz1} in the finite case) 
we see that $\bVi(\Tn{\bbi})\cong {}_iP$. As discussed 
in Subsection~\ref{sub-ammod}, the ${}_iP$ can be given 
a grading coming from the path length $l(\cdot)$. 
Thus, all the $\Tn{\bbi}$'s are gradable.




On the walls: Since this is the semisimple 
case by Lemma~\ref{lem-block}, we can just 
assign to any simple $\bVi(\Tn{\mu_i})$ a 
degree by demanding that all elements of 
$\bVi(\Tn{\mu_i})$ are concentrated in 
this particular degree. Thus, all 
$\Tn{\mu_i}$ are gradable and concentrated in (up to shifts) degree zero.



The uniqueness (up to shifts) can be proven as in~\cite[Lemma~2.5.3]{bgs}.
\end{proof}



\begin{cor}\label{cor-gradings2}
We have isomorphisms of ungraded categories 
$\T_{\lambda}(\leq\bbm)\cong_{\mathrm{iso}}\Tgr(\leq\bbm)$ 
and $\T\cong_{\mathrm{iso}}\Tgr$. That is, 
also all morphisms are gradable and the gradings are unique up to a shift.\makeqed
\end{cor}



\begin{proof}
As before, for any 
two gradable $A$-homomorphisms $f$ and 
$f^{\prime}$, the direct sum $f\oplus f^{\prime}$ 
is also gradable. Thus, by Lemma~\ref{lem-tiltcat} 
part (a) again together with Proposition~\ref{prop-grading}, 
it is enough to consider only 
$\Hom_{\Uq}(\Tn{\bbi},\Tn{\bbi^{\prime}})$. 
By Corollary~\ref{cor-homspaces} we see that 
this space has a basis consisting of a subset 
of $\bbi,u_i,d_i$ and $\varepsilon_i$ (in the non-semisimple case) who are all gradable by Proposition~\ref{prop-isoofalg2}. The semisimple 
case and the uniqueness of the grading follow as 
above. This gives rise to an 
isomorphism after collapsing the grading of $\Tgr$ 
(since each object $M$ of $\T$ corresponds to objects of $\Tgr$ 
of the form $M\langle s\rangle$ for $s\in\bZ$).
\end{proof}



Using Corollary~\ref{cor-gradings2} and 
abuse of language, we do not distinguish 
between $\Uq$-tilting modules or their graded versions.
Moreover, as a consequence of 
Corollary~\ref{cor-gradings2}, we can 
choose a \textit{standard grading} 
(since it will be unique up to shifts) 
by demanding that simple $\Uq$-tilting modules 
are concentrated in degree zero and 
proceed inductively along the quiver. 
Note that, by Propositions~\ref{prop-weyl} 
and~\ref{prop-tilt} and Corollary~\ref{cor-weyl}, 
this gives inductively a graded lift of 
$\Ln{\bbi}$, $\Vn{\bbi}$ and $\dVn{\bbi}$ as well. 
This choice gives rise to the following graded 
refinements of Propositions~\ref{prop-weyl} 
and~\ref{prop-tilt} and Corollary~\ref{cor-weyl}.



\begin{prop}\label{prop-refine}
Suppose $i=al+b$ for some $a,b\in\bN$ with 
$b\leq l-2$. Set $i^{\prime}=(a+2)l-b-2$. 
Then there
exist short exact sequences
\begin{gather*}
0\longrightarrow \Ln{i}\langle+1\rangle\hooklongrightarrow\Vn{i^{\prime}}\twoheadlongrightarrow \Ln{i^{\prime}}\longrightarrow 0,
\hspace*{0.25cm}
0\longrightarrow \Ln{i^{\prime}}\hooklongrightarrow\dVn{i^{\prime}}\twoheadlongrightarrow \Ln{i}\langle-1\rangle\longrightarrow 0,
\\
0\longrightarrow \Vn{i^{\prime}}\hooklongrightarrow\Tn{i^{\prime}}\twoheadlongrightarrow \Vn{i}\langle-1\rangle\longrightarrow 0,
\hspace*{0.25cm}
0\longrightarrow \dVn{i}\langle+1\rangle\hooklongrightarrow\Tn{i^{\prime}}\twoheadlongrightarrow \dVn{i^{\prime}}\longrightarrow 0.
\end{gather*}
Moreover, all involved morphisms are degree preserving.\makeqed
\end{prop}



\begin{proof}
This follows from our choice for the grading convention 
and the degree (under the equivalence in 
Theorems~\ref{thm-struktursatz1} and~\ref{thm-struktursatz2}) 
of the morphisms in $\Hom_{\Uq}(\Tn{\bbi},\Tn{\bbi^{\prime}})$. 
To be more precise, we know that the $\Tn{\bbi}$'s will 
be, under the functor $\bVi$, 
mapped to ${}_iP$. Then for example, 
as explained in Proposition~\ref{prop-link}, the unique morphism
\[
\Tn{\bbi}\twoheadrightarrow\dVn{\bbi}\hookrightarrow\Tn{\bbiu}\leftrightsquigarrow\xymatrix{
  \raisebox{0.1cm}{\xy(0,0)*{\bullet};(0,-2.5)*{\scriptstyle \bbiu};\endxy}\!\!   &  \raisebox{0.1cm}{\xy(0,0)*{\bullet};(0,-2.5)*{\scriptstyle \bbi};\endxy} \ar[l]_/-0.2em/{u_i}\\
}
\]
is of degree $1$. By our convention 
we see that moving to the left along 
the KS $\infty$-quiver always increases 
the degree (starting in degree zero for the first simple module). Thus,
\[
0\longrightarrow \dVn{i}\langle+1\rangle\hooklongrightarrow\Tn{i^{\prime}}\twoheadlongrightarrow \dVn{i^{\prime}}\longrightarrow 0.
\]
All other cases follow similarly 
and are left to the reader. Note again 
that moving to the left along the KS 
$\infty$-quiver always increases the degree, 
but, by duality, moving right decreases it.
\end{proof}



\begin{rem}\label{rem-webalg2}
As we already mentioned in 
Remark~\ref{rem-webalg} the grading induced on 
$\T$ comes from an Euler characteristic on a 
certain cobordism category associated to 
Khovanov's arc algebra $H^m$. Thus, this is 
a ``natural'' grading from the viewpoint of topology.



As we mentioned above, we think that this 
should generalize in type $\boldsymbol{\mathrm{A}}$. The degree will 
there be given by an Euler characteristic on 
a certain ``foam'' category associated to the 
$\mathfrak{sl}_n$-web algebras which generalize Khovanov's arc algebra. See for example~\cite{lqr1},~\cite{mpt},~\cite{qr1} or~\cite{tub4}. 
Or alternatively, by a grading coming from certain cyclotomic KL-R algebras in the sense of~\cite{kl5} or~\cite{rou}.
\end{rem}




We conclude this section by lifting the translation 
onto $\mathcal T^{\mu}_{\lambda}$, the translation out 
of $\mathcal T_{\mu}^{\lambda}$ and the 
translation through the wall $\Theta_s$ and $\Theta_t$ 
to their graded versions (we use the same notation for these). 
To understand our notation, recall that a category 
$\tilde{\mathcal C}$ is called a \textit{graded} 
category, if there exists a functor 
$\mathrm{deg}\colon\tilde{\mathcal C}\to\mathcal Z$ 
where we consider $\bZ$ as a category $\mathcal Z$ 
with $\Ob(\mathcal Z)=\{\bullet\}$ and 
$\End_{\mathcal Z}(\bullet)=\bZ$. In addition, there is a 
$\bZ$-action $\langle s\rangle$ on the objects 
$O\in\Ob(\tilde{\mathcal C})$, called \textit{shift} 
by $s\in\bZ$, such that a given morphism 
$f\in\Hom_{\tilde{\mathcal C}}(O,O^{\prime})$ of 
degree $d$ is of degree $d+s^{\prime}-s$ in 
$\Hom_{\tilde{\mathcal C}}(O\langle s\rangle,O^{\prime}\langle s^{\prime}\rangle)$ 
(said otherwise, 
every morphism in $\tilde{\mathcal C}$ has a degree which behaves 
additively under composition and objects can be shifted via 
$\langle s\rangle$). Moreover, denote by $\mathcal C$ the 
\textit{ungraded version}. We say a functor 
$\mathcal F\colon\mathcal C\to\mathcal D$ is gradable 
if there exists a \textit{lift} 
$\tilde{\mathcal F}\colon\tilde{\mathcal C}\to\tilde{\mathcal D}$ such that
$\mathrm{forget}\circ\tilde{\mathcal F}=\mathcal F\circ\mathrm{forget}$.
We say that $\tilde{\mathcal F}$ 
and $\tilde{\mathcal F}^{\prime}$ are the 
\textit{same up to a shift $s\in\bZ$} if there 
is a natural isomorphism\footnote{In the $2$-category 
where objects are graded categories, where morphisms 
are graded functors between these categories and 
where $2$-morphisms are graded natural transformations.} 
between $\tilde{\mathcal F}$ and $\tilde{\mathcal F}^{\prime}\langle s\rangle$.
As before, we are usually very 
careless with our distinction 
of $\tilde{\mathcal F}$ and $\mathcal F$.


\begin{cor}\label{cor-funcgrad2}
The functors 
$\mathcal T_{\lambda}^{\mu,i},\mathcal T_{\mu}^{\lambda,i},\mathcal T_{\lambda}^{\mu},\mathcal T_{\mu}^{\lambda}\in\Ob(\Endo(\T(\leq\bbm)))$
are gradable (same in the infinite case). 
Moreover, any two lifts of 
$\mathcal T_{\lambda}^{\mu,i},\mathcal T^{\lambda,i}_{\mu}$ are the same up to shifts.\makeqed
\end{cor}



\begin{proof}
This follows by Remark~\ref{rem-othermaps} 
and the fact that $\mathcal U_{i}^t$ and $\mathcal U^i_t$ are 
graded functors (similarly in the $\infty$ case). 
The statement for 
$\mathcal T_{\lambda}^{\mu}$ and $\mathcal T_{\mu}^{\lambda}$ follows now 
by~\eqref{eq-projectivefunc2}, since clearly sums of gradable, 
additive functors in an additive category are 
gradable. For the infinite case note that any sum (finite or not) 
of gradable functors is gradable and the claim 
follows from~\eqref{eq-projectivefunc2} 
and the finite case.
The uniqueness up to shifts can again be verified 
as in~\cite[Lemma 2.5.3]{bgs}, since 
$\mathcal T_{\lambda}^{\mu,i}$ and 
$\mathcal T_{\mu}^{\lambda,i}$ are 
indecomposable by Proposition~\ref{prop-profunc}. 
\end{proof}



\begin{cor}\label{cor-funcgrad1}
The functors 
$\Theta^{\lambda,i}_s,\Theta^{\lambda,i}_t,\Theta^{\lambda}_s,\Theta^{\lambda}_t\in\Ob(\Endo(\T_{\lambda}(\leq\bbm)))$ 
are 
gradable (same in the infinite case). 
Moreover, any two lifts of $\Theta^{\lambda,i}_{s\text{ or }t}$ are the same up to a 
shift.\makeqed
\end{cor}



\begin{proof}
We can use Corollary~\ref{cor-funcgrad2} 
since all 
functors involved can be obtained via sums and compositions of the functors 
$\mathcal T_{\lambda,i}^{\mu}$ and $\mathcal T^{\lambda,i}_{\mu}$. The infinite case follows as above, but using~\eqref{eq-projectivefunc} 
instead of~\eqref{eq-projectivefunc2}. 
By Proposition~\ref{lem-profuncinde} the functors $\Theta^{\lambda,i}_{s\text{ or }t}$ are 
indecomposable and the uniqueness up to shifts follows as before. 
\end{proof}


We choose the graded versions of 
$\mathcal T_{\lambda}^{\mu,i}$ and out of $\mathcal T^{\lambda,i}_{\mu}$ 
in the evident way as induced by the gradings on $\mathcal U_i^t$ or 
$\mathcal U^i_t$.
We get the following refinement of Proposition~\ref{prop-functors}.
\begin{cor}\label{cor-funcgrad3}
In the graded category $\Tgr$ we have for all 
$i\in\bN$ the following.
\begin{itemize}
\item[(a)] The functors 
$\mathcal T_{\lambda}^{\mu}$ and $\mathcal T^{\lambda}_{\mu}$ are 
well-defined (their definition gives gradable $\Uq$-tilting modules in the 
right blocks), (up to shifts) 
adjoints (left and right) and 
exact. Thus, $\Theta^{\lambda}_s$ and $\Theta^{\lambda}_t$ are exact and graded self-adjoint.
\item[(b)] We have (recalling that $\Ln{\mu_i}\cong\Tn{\mu_i}$)
\begin{gather*}
\mathcal T_{\lambda}^{\mu}(\Tn{\lambda_i})\cong\begin{cases}\Tn{\mu_{i-1}}\langle+1\rangle\oplus\Tn{\mu_{i+1}}\langle+1\rangle,  & \text{if }\mu_i>\lambda_i,\\
  \Tn{\mu_i}\oplus\Tn{\mu_i}\langle+2\rangle,  & \text{if }\mu_i<\lambda_i,\end{cases}
\\
\mathcal T^{\lambda}_{\mu}(\Tn{\mu_i})\cong\begin{cases}\Tn{\lambda_{i+1}}\langle-1\rangle,  & \text{if }\mu_i>\lambda_i,\\
 \Tn{\lambda_{i}}\langle-1\rangle,  & \text{if }\mu_i<\lambda_i.\end{cases}
\end{gather*}
\item[(c)] The dead-end relations $\Theta^{\lambda}_s(\Tn{\lambda_0})\cong0$, $\Theta^{\lambda}_s(\Tn{\lambda_1})\cong\Tn{\lambda_2}$ and $\Theta^{\lambda}_t(\Tn{\lambda_0})\cong\Tn{\lambda_1}$. Moreover, we have
\[
\Theta^{\lambda}_{s\text{ or }t}(\Tn{\lambda_i})\cong\begin{cases}\Tn{\lambda_{i-1}}\oplus\Tn{\lambda_{i+1}} ,  & \text{if }i>1\text{ is odd for }s\text{ and even for }t,\\
  \Tn{\lambda_{i}}\langle-1\rangle\oplus\Tn{\lambda_{i}}\langle+1\rangle ,  & \text{if }i>0\text{ is odd for }t\text{ and even for }s.\end{cases}
\]
\end{itemize}
Here we set $\Tn{-1}=\Tn{\lambda_{-1}}=\Tn{\mu_{-1}}=0$.\makeqed
\end{cor}



\begin{proof}
This follows directly from 
Theorem~\ref{thm-struktursatz2} 
parts (d), (f) and (g) together 
with~\eqref{eq-ufunctors} and the 
discussion in Subsection~\ref{sub-endo}, for example Lemma~\ref{lem-adjoint}.
\end{proof}



\begin{prop}\label{prop-endograd}
All functors in $\Endo(\T)$ are gradable. 
Moreover, the grading is unique up to shifts 
on the indecomposable projective functors. 
Similar for $\Endo(\T_{\lambda})$.\makeqed
\end{prop}



\begin{proof}
The uniqueness, up to shifts, of the grading on the indecomposable 
factors can be proven as before. By the Krull-Schmidt property of $\Endo(\T)$ 
(see Lemma~\ref{lem-profunc}) and the fact that any direct sum of gradable, 
additive functors in an additive category is 
gradable, it suffices to show that indecomposable 
projective functors are gradable. By 
Lemma~\ref{lem-profuncinde} these are given 
by tensoring with a $\Tn i$ for some $i\in\bN$. 
By Proposition~\ref{prop-grading} all objects 
of $\T$ are gradable and, by Lemma~\ref{lem-tiltcat} 
part (b), the category $\T$ is preserved by finite 
tensor products. Thus, the indecomposable projective functors are gradable. 
For 
$\Endo(\T_{\lambda})$ this follows from Corollary~\ref{cor-funcgrad1}.
\end{proof}



Proposition~\ref{prop-endograd} motivates 
the definition of $\Endo(\Tgrl)$ which can be 
defined in the spirit of~\eqref{eq-gradhom} (thus, using 
\textit{graded} hom-spaces of natural 
transformations). This is the same as 
saying that the category $\Endo(\Tgrl)$ is graded in the above sense.



In order to state the proposition, we write 
$\bVi\mathcal F(\Tn{\bbin})$ short for 
$\bigoplus_i\bVi\mathcal F(\Tn{\bbi})$. This 
is an $\Ai$-bimodule via pre-composition 
(right) and post-composition (left) with 
$\mathcal F(\cdot)$ or $\mathcal F^{(\prime)}(\cdot)$.



Moreover, recall that the superscript fs 
indicates that we are only considering finitely 
supported homomorphisms. We have the following 
(similar of course for $\Endo(\Tgr_{\lambda}(\leq\bbm))$).



\begin{prop}\label{prop-endograd2}
Let $\mathcal F,\mathcal F^{\prime}$ be two 
functors in $\Endo(\Tgrl)$. Then there exists 
an isomorphism of graded $\bC$-vector spaces
\[
\Hom^{\mathrm{fs}}_{\Endo(\Tgrl)}(\mathcal F,\mathcal F^{\prime})\cong \Hom^{\mathrm{fs}}_{\AipModgrAi}(\bVi\mathcal F(\Tn{\bbin}),\bVi\mathcal F^{\prime}(\Tn{\bbin})).
\]
This is an isomorphism of graded rings if $\mathcal F=\mathcal F^{\prime}$.\makeqed
\end{prop}



\begin{proof}
It 
suffices to verify the statement for the 
translation onto $\mathcal T^{\mu}_{\lambda}$, 
out of $\mathcal T_{\mu}^{\lambda}$ and the 
through the wall functors $\Theta^{\lambda}_{s\text{ or }t}$. 
By Theorem~\ref{thm-struktursatz2} part (d), (f) and (g) 
it suffices to show the statement for the corresponding $\mathcal U$'s.



Without loss of generality, we only discuss the 
``even'' case, that is, $\mathcal U^t_{\mathrm{even}}$, $\mathcal U_t^{\mathrm{even}}$ 
and $\mathcal U^{\infty}_{\mathrm{even}}$ which correspond to 
$\mathcal T^{-1}_{\lambda}$, $\mathcal T_{-1}^{\lambda}$ and 
$\Theta^{\lambda}_{s}$ respectively. Thus, by additivity, it 
suffices to show the statement for $\mathcal U^t_{i}$, 
$\mathcal U_t^{i}$ and $\mathcal U_i$ for $i>0$ even.
Fix $i$ and assume that $i\ll m$. Recall 
that Khovanov and Seidel have shown 
in~\cite[Proposition~2.9]{ks1} that 
$\pModAm$ and $\mathcal O^{\mathfrak{p}}_0$ 
for $\mathfrak{sl}_{m+1}$ and parabolic 
$\mathfrak{p}$ of type $S_1\times S_{m-1}$ 
are equivalent. Moreover, as they explain 
below~\cite[Proposition~2.9]{ks1}, 
$\mathcal U^t_{i}$, $\mathcal U_t^{i}$ and 
$\mathcal U_i$ correspond to the translation onto, out off 
and through the $i$-th wall functors.



Note now that the same holds for the graded 
version of $\mathcal O^{\mathfrak{p}}_0$ 
introduced in~\cite{st2} as Stroppel explains in~\cite[Section~1]{st1}. 
Thus, we can use~\cite[Theorem~1.12]{st1} to finish the proof in the 
``cut-off case''. Taking the inductive limit now shows the statement.
\end{proof}



\begin{cor}\label{cor-center}
There is an isomorphism of graded rings
\[
\End^{\mathrm{fs}}_{\mathrm{gr}}(\mathrm{id})\cong Z(\Ai),
\]
where $\mathrm{id}$ is the identity functor on 
$\Tgrl$, $Z(\Ai)$ is the center of 
$\Ai$ and the endomorphism ring is to be taken in $\Endo(\Tgrl)$.\makeqed
\end{cor}



\begin{proof}
From Proposition~\ref{prop-isoofalg2} 
combined with Proposition~\ref{prop-endograd2}. 
Note that we get the center $Z(\Ai)$ and not $\Ai$ 
itself because the endomorphism ring from 
Proposition~\ref{prop-isoofalg2} is as 
$\Uq$-intertwiners while the one induced from 
Proposition~\ref{prop-endograd2} is as $\AipModgrAi$-bimodule intertwiners.
\end{proof}



\begin{rem}\label{rem-grothendieck1a}
It is known 
that the split Grothendieck group 
$K_0^{\oplus}(\cdot)=K_0^{\oplus}(\cdot)\otimes_{\bZ[v,v^{-1}]}\bC$ 
(viewed as a module over $\bZ[v,v^{-1}]$ where the 
formal parameter $v$ comes for the grading) of 
$\pModgrAm$ \textit{categorifies} the Burau 
representation $\mathcal B_{m+1}$ of the 
$m+1$-strand braid group $B_{m+1}$ 
(the action of $B_{m+1}$ is induced via 
functors constructed from $\mathcal U_i$), see~\cite[Proposition~2.8 and Subsection 2.d]{ks1}. 
Using Theorem~\ref{thm-struktursatz1}, we see that the 
same is true for $\Tgrl(\leq\bbm)$, that is 
$K_0^{\oplus}(\Tgrl(\leq\bbm))\cong\mathcal B_{m+1}$ as 
$B_{m+1}$-modules. Note now that the limit therefore 
categorifies the corresponding $\infty$ version of 
the Burau representation of the braid 
group $B_{\infty}$ with $\infty$-many strands.
\end{rem}



\begin{rem}\label{rem-grothendieck1}
Fix $\mathfrak{sl}_{m+1}$ and denote 
by $\mathfrak{p}_i$ the parabolic for 
$S_i\times S_{m-i}$. Moreover, we denote 
by $\mathcal O^{\mathfrak{p}}_{\mathrm{max}}=\bigoplus_{i=0}^m\mathcal O^{\mathfrak{p}_i}_{0}$ 
the direct sum. We note that, as Bernstein-Frenkel-Khovanov 
conjectured in~\cite[Conjectures~1-4]{bfk} and Stroppel proved 
in~\cite{st3}, parabolic category 
$\mathcal O^{\mathfrak{p}}_{\mathrm{max}}$ 
for $\mathfrak{sl}_{m+1}$ (compare to Remark~\ref{rem-cato}) 
can be used to \textit{categorify} the $m+1$-strand 
Temperley-Lieb algebra $TL^v_{m+1}$. The split 
Grothendieck group $K_0^{\oplus}$ of the category 
of projective endofunctors on 
$\mathcal O^{\mathfrak{p}}_{\mathrm{max}}$ gives\footnote{The parameter $q$ in the notation of Bernstein-Frenkel-Khovanov comes from the 
grading.} $TL^v_{m+1}$. 
The Grothendieck group behaves additive, i.e.
\[
TL^v_{m+1}\cong K_0^{\oplus}(\Endo(\mathcal O^{\mathfrak{p}}_{\mathrm{max}}))\cong \bigoplus_{i=0}^mK_0^{\oplus}(\Endo(\mathcal O^{\mathfrak{p}}_{i})).
\]
Thus, $\Endo(\mathcal O^{\mathfrak{p}}_{1})$ 
gives a summand $\overline{TL}^v_{m+1}$ of $TL^v_{m+1}$ 
that corresponds to the ``next to highest weight'' 
summand of $V^{\otimes m}$ (with notation as in 
Remark~\ref{rem-webalg}), see~\cite[Section~6]{chkh}.



Thus, as explained in Remark~\ref{rem-cato}, 
by Proposition~\ref{prop-endograd2}, the same 
is true for $\Endo(\Tgrl(\leq\bbm))$. Moreover, 
it is easy to see that the embedding of categories
\[
\Endo(\Tgrl(\leq\bbm))\hookrightarrow\Endo(\Tgrl(\leq\bbmu))
\]
gives rise to an embedding
\[
\overline{TL}^v_{m+1}\cong K_0^{\oplus}(\Endo(\Tgrl(\leq\bbm)))\hookrightarrow K_0^{\oplus}(\Endo(\Tgrl(\leq\bbmu)))\cong \overline{TL}^v_{m+2}.
\]
Thus, by Proposition~\ref{prop-isoofalg2}, $\Endo(\Tgrl)$ 
governs all of them at once.
\end{rem}